% ===========================================================================
%  section-double-end.tex
%  "A two-anchor normal form" -- drop-in section for main.tex.
%
%  ASSUMPTIONS ON THE PREAMBLE.  This file introduces no packages and no
%  macros of its own.  It assumes exactly the following, all of which are
%  already provided by the preamble of main.tex:
%
%    packages   amsmath, amssymb, amsthm  (align, \binom, \sqcup, \operatorname)
%               booktabs                  (\toprule \midrule \bottomrule)
%               xcolor                    (used only through \todo)
%               natbib                    (\cite)
%    theorem environments (plain style, shared counter, numbered within
%               section):  theorem, lemma, proposition, corollary
%               (remark style): remark, remark*
%    macros     \bin{a}{b}   = \binom{a}{b}
%               \diam        = \operatorname{diam}
%               \Z           = \mathbb{Z}
%               \file{...}, \tok{...}   (typewriter)
%               \todo{...}   (visible red marker; \todofalse hides them)
%               \bk          (discretionary break; used inside \tok only)
%
%    labels referenced from elsewhere in main.tex (all exist there):
%               sec:prelim, sec:algorithm, sec:verification, sec:crowding,
%               sec:endpoint, sec:data, sec:ai,
%               lem:taylor, lem:golomb, thm:endpoint, thm:crowd280,
%               thm:crowd283, cor:reduction, tab:known, tab:crowdhash
%    bib keys:  Leech1975, Taylor1977, SzekelyWangZhang2005, Calhoun2007,
%               OEIS, Shearer1990, Dollas1998, DistributedNetOGR
%
%  Notation follows main.tex: n is the order, N = \bin n2, and e, f, p, h are
%  the quantities already introduced inside the proof of Theorem~\ref{thm:endpoint}.
% ===========================================================================

\section{A two-anchor normal form}\label{sec:doubleend}

The engine of Section~\ref{sec:algorithm} is organised by edges, and the
order-$25$ argument of Section~\ref{sec:crowding} is organised by a deleted
leaf and carries a parameter $\delta=\diam(T-\ell)$.  This section develops a
third normal form, organised instead by the two endpoints of the unique
diametral pair.  It has two features that the other two lack.  First, it has no
parameter: the whole order-$n$ problem becomes one finite additive covering
problem, stated once.  Second, the objects it enumerates are not trees but
pairs of integer sets, so the search order can be driven by the target values
$1,2,\dots,N$ rather than by the shape of the tree.

We state the normal form (Section~\ref{sec:de:setup}), derive from it the exact
parity split, an attachment lemma, a Golomb bound on the spine and a ball bound
(Section~\ref{sec:de:cons}), record a generating-function identity together
with an honest account of how little its most quotable corollary is worth
(Section~\ref{sec:de:lemmad}), describe the search built on it
(Section~\ref{sec:de:search}) and its validation
(Section~\ref{sec:de:valid}), and close with a statement of exactly what has
and has not been computed (Section~\ref{sec:de:status}).  The engine is
independent of every other program in this paper: it shares no code with the
forest search, with the per-topology search or with the order-$25$ material of
Section~\ref{sec:crowding}.

\subsection{Coordinates at the two diametral endpoints}\label{sec:de:setup}

Throughout this section $T$ is a Leech tree of order $n\ge3$, $N=\bin n2$ and
$V=n-2$.  As recorded at the start of Section~\ref{sec:endpoint}, the value $N$
is attained by exactly one pair $\{a,b\}$ and both $a$ and $b$ are leaves; we
call $\{a,b\}$ the \emph{anchors} and write $P$ for the $a$--$b$ path, of weight
$N$.  For a vertex $u\notin\{a,b\}$ let $q(u)$ be the vertex of $P$ nearest to
$u$ and put
\[
  p(u)=d\bigl(a,q(u)\bigr),\qquad h(u)=d\bigl(u,q(u)\bigr)\ge0,
\]
so that $d(a,u)=p(u)+h(u)$ and $d(b,u)=\bigl(N-p(u)\bigr)+h(u)$.  These are the
quantities used in the proof of Theorem~\ref{thm:endpoint}; here they become
standing notation.

For any two vertices $x\ne y$ of $T$ we call
$\operatorname{gap}(x,y)=N-d(x,y)$ the \emph{gap} of the pair $\{x,y\}$, and for
$u\notin\{a,b\}$ we write
\[
  e(u)=\operatorname{gap}(a,u)=N-d(a,u),\qquad
  f(u)=\operatorname{gap}(b,u)=N-d(b,u),
\]
the \emph{gap coordinates} of $u$, and set $E=\{e(u):u\notin\{a,b\}\}$ and
$F=\{f(u):u\notin\{a,b\}\}$.

The two coordinates determine the position of $u$:
\begin{equation}\label{eq:deph}
  p(u)=\frac{N+f(u)-e(u)}2,\qquad h(u)=\frac{N-e(u)-f(u)}2 .
\end{equation}
The Leech condition says that the $N$ pairs of $T$ have distinct distances
filling $[1,N]$, i.e.\ that their $N$ gaps are exactly $0,1,\dots,N-1$; the
pair $\{a,b\}$ contributes the gap $0$, the pairs $\{a,u\}$ and $\{b,u\}$
contribute $E$ and $F$, and the remaining $\bin V2$ pairs contribute the values
computed in Lemma~\ref{lem:degap} below.

\begin{lemma}[coordinate constraints]\label{lem:decoord}
For every $u\notin\{a,b\}$,
\begin{enumerate}
\item[(i)] $1\le e(u)\le N-1$ and $1\le f(u)\le N-1$;
\item[(ii)] $e(u)+f(u)=N-2h(u)\le N$, with equality if and only if $u$ lies on $P$;
\item[(iii)] $e(u)+f(u)\equiv N\pmod 2$;
\item[(iv)] $E$ and $F$ each have exactly $V$ elements and $E\cap F=\varnothing$;
in particular $e(u)\ne f(u)$, so no vertex sits at $p=N/2$.
\end{enumerate}
\end{lemma}

\begin{proof}
$d(a,u)<N$ because $\{a,b\}$ is the only pair at distance $N$ and $u\ne b$, and
$d(a,u)\ge1$, which gives (i) for $e$; symmetrically for $f$.  Adding
$d(a,u)=p+h$ and $d(b,u)=(N-p)+h$ gives $d(a,u)+d(b,u)=N+2h(u)$, whence (ii)
and (iii).  For (iv), the pairs $\{a,u\}$ are pairwise distinct, so the
distances $d(a,u)$ are pairwise distinct and $|E|=V$; likewise $|F|=V$; and
$\{a,u\}\ne\{b,v\}$ for all $u,v\notin\{a,b\}$, so $d(a,u)\ne d(b,v)$ and
$E\cap F=\varnothing$.  Taking $u=v$ gives $e(u)\ne f(u)$, which by
\eqref{eq:deph} says $p(u)\ne N/2$.
\end{proof}

\begin{remark*}
The parity statement is $e+f\equiv N$, not $e\equiv f$.  The two agree only
when $N$ is even, that is for $n\equiv0,1\pmod4$; for the Leech tree of order
$3$ ($N=3$) one has $E=\{2\}$, $F=\{1\}$.  Both cases occur among the known
Leech trees and both are used below, so we keep them apart.
\end{remark*}

\begin{lemma}[crossing formula]\label{lem:degap}
Let $u\ne v$ be vertices outside $\{a,b\}$ and let $h_c(u,v)\ge0$ be the height
above $P$ of the last common vertex of the $u$--$P$ and $v$--$P$ paths.  Then
\begin{equation}\label{eq:degap}
  \operatorname{gap}(u,v)\;=\;\min\bigl(f(u)+e(v),\;f(v)+e(u)\bigr)\;+\;2\,h_c(u,v),
\end{equation}
where $h_c(u,v)=0$ unless $f(u)-e(u)=f(v)-e(v)$, equivalently $p(u)=p(v)$, in
which case $0\le h_c(u,v)\le\min\bigl(h(u),h(v)\bigr)$.  Moreover the two
\emph{crossings} $f(u)+e(v)$ and $f(v)+e(u)$ differ by $2\bigl(p(v)-p(u)\bigr)$
and in particular always have the same parity.
\end{lemma}

\begin{proof}
By \eqref{eq:deph}, $f(u)-e(u)=2p(u)-N$, so
$\bigl(f(v)+e(u)\bigr)-\bigl(f(u)+e(v)\bigr)=2\bigl(p(v)-p(u)\bigr)$; hence
$p(u)\le p(v)$ if and only if $f(u)+e(v)\le f(v)+e(u)$, and the two crossings
are equal exactly when $p(u)=p(v)$.  Assume $p(u)\le p(v)$.  If
$q(u)\ne q(v)$, the $u$--$v$ path leaves $u$, meets $P$ at $q(u)$, runs along
$P$ to $q(v)$ and leaves $P$ to $v$, so
$d(u,v)=h(u)+\bigl(p(v)-p(u)\bigr)+h(v)=N-f(u)-e(v)$, and $h_c=0$; the same
computation applies when $q(u)=q(v)$ but the two paths leave that vertex by
different edges.  Otherwise $q(u)=q(v)$ and the two paths share an initial
segment ending at height $h_c\ge0$, and
$d(u,v)=h(u)+h(v)-2h_c=N-f(u)-e(v)-2h_c$.  Taking gaps gives \eqref{eq:degap};
$h_c\le\min(h(u),h(v))$ because the common vertex lies on both paths.
\end{proof}

\begin{theorem}[two-anchor normal form]\label{thm:denormal}
Let $T$ be a weighted tree of order $n\ge3$ and let $a,b$ be two of its vertices
with $d(a,b)=N=\bin n2$; define $p,h,e,f$ relative to the $a$--$b$ path as
above.  Then $T$ is a Leech tree if and only if
\begin{equation}\label{eq:detiling}
  \{0\}\ \sqcup\ E\ \sqcup\ F\ \sqcup\
  \bigl\{\operatorname{gap}(u,v):\{u,v\}\subseteq V(T)\setminus\{a,b\}\bigr\}
  \;=\;[0,N-1],
\end{equation}
the union being disjoint and the
$1+V+V+\bin V2=N$ values pairwise distinct.  In that case $\{a,b\}$ is the
unique diametral pair, $a$ and $b$ are leaves, and Lemma~\ref{lem:decoord}
holds.
\end{theorem}

\begin{proof}
The $N$ pairs of $T$ split into $\{a,b\}$, the $V$ pairs $\{a,u\}$, the $V$
pairs $\{b,u\}$ and the $\bin V2$ pairs $\{u,v\}$ with $u,v\notin\{a,b\}$, and
their gaps are respectively $0$, the elements of $E$, the elements of $F$ and
the values \eqref{eq:degap}, Lemma~\ref{lem:degap} being valid in any weighted
tree.  The Leech condition is that the $N$ distances are distinct and fill
$[1,N]$, i.e.\ that the $N$ gaps are distinct and fill $[0,N-1]$, which is
\eqref{eq:detiling}.  Finally
$1+2V+\bin V2=1+2(n-2)+\bin{n-2}2=\bin n2=N$, and the last sentence is the
first paragraph of Section~\ref{sec:endpoint} together with
Lemma~\ref{lem:decoord}.
\end{proof}

Theorem~\ref{thm:denormal} is a reformulation, not a new obstruction; its value
is that the right-hand side is a fixed interval, the left-hand side is built
from sums of the unknowns, and no parameter such as $\delta$ appears.  A single
search over the pairs $(e_i,f_i)$ and the tie heights $h_c$ therefore covers
the whole order-$n$ problem.

\begin{remark}[what the relaxation is, and is not]\label{rem:derelax}
The system consisting of Lemma~\ref{lem:decoord}, formula~\eqref{eq:degap} with
$0\le h_c\le\min(h_u,h_v)$ and $h_c=0$ for untied pairs, and the
partition~\eqref{eq:detiling} is a set of \emph{necessary} conditions on an
abstract pair of coordinate systems.  It is strictly weaker than being a tree:
solutions of the relaxation need not be realisable, because nothing in it
forces the branch carrying a vertex with $h>0$ to hang on a vertex of $T$.  At
$n=9$ the search with every condition enforced \emph{except} the attachment
condition of Lemma~\ref{lem:deattach} returns exactly $18$ solutions, and none
of them is a tree (Section~\ref{sec:de:valid}).  Consequently a search over
this system proves non-existence when it is exhausted and empty, while every
solution it does report must be converted into an explicit weighted tree and
re-checked before it is believed.  Both directions are implemented and used.
\end{remark}

\subsection{Immediate consequences}\label{sec:de:cons}

\subsubsection*{The parity split}

\begin{proposition}[exact parity split]\label{prop:desplit}
Let $T$ be a Leech tree of order $n\ge3$.  Colour a non-anchor vertex $u$
\emph{black} if $e(u)$ is even and \emph{white} if $e(u)$ is odd, and let
$\alpha$ and $\beta$ be the numbers of black and of white non-anchor vertices,
so that $\alpha+\beta=V$.  Extend the colouring to the anchors as follows:
if $N$ is even, colour both $a$ and $b$ black; if $N$ is odd, colour $b$ black
and $a$ white.  Call a pair of vertices \emph{bichromatic} if its two endpoints
have different colours.  Then
\begin{enumerate}
\item[(a)] if $N$ is even, $e(u)\equiv f(u)\pmod2$ for every $u$, a pair has an
odd gap if and only if it is bichromatic, and counting bichromatic pairs gives
\[
  (\alpha+2)\,\beta=\tfrac N2=\tfrac{n(n-1)}4,
  \qquad\text{hence}\qquad
  \beta=\frac{n\pm\sqrt n}2,\quad \alpha=n-2-\beta ;
\]
\item[(b)] if $N$ is odd, $e(u)\not\equiv f(u)\pmod2$ for every $u$, a pair has
an \emph{even} gap if and only if it is bichromatic, and counting bichromatic
pairs gives
\[
  (1+\alpha)(1+\beta)=\Bigl\lceil\tfrac N2\Bigr\rceil=\tfrac{N+1}2,
  \qquad\text{hence}\qquad
  \{\alpha,\beta\}=\Bigl\{\tfrac{(n-2)-\sqrt{n-2}}2,\ \tfrac{(n-2)+\sqrt{n-2}}2\Bigr\}.
\]
\end{enumerate}
In both cases $0\le\alpha,\beta\le V$, which can exclude one of the two roots.
\end{proposition}

\begin{proof}
By Lemma~\ref{lem:degap} the two crossings of a pair have the same parity, so
the parity of $\operatorname{gap}(u,v)$ is that of $f(u)+e(v)$, the correction
$2h_c$ being even.  Write $\varepsilon(u)\in\{0,1\}$ for the parity of $e(u)$,
so that black means $\varepsilon=0$; by Lemma~\ref{lem:decoord}(iii) the parity
of $f(u)$ is $\varepsilon(u)+N$.

Suppose $N$ is even, so $f(u)\equiv e(u)$ and both anchors are black.  The pair
$\{a,b\}$ has gap $0$ and is monochromatic.  The pair $\{a,u\}$ has gap $e(u)$
and is bichromatic exactly when $u$ is white, i.e.\ exactly when $e(u)$ is odd;
the pair $\{b,u\}$ has gap $f(u)\equiv e(u)$ and the same holds.  Finally
$\operatorname{gap}(u,v)\equiv f(u)+e(v)\equiv\varepsilon(u)+\varepsilon(v)$, so
it is odd exactly when $u$ and $v$ have different colours.  Hence the odd gaps
are precisely the bichromatic pairs, of which there are $(\alpha+2)\beta$, and
$[0,N-1]$ contains $N/2$ odd values.  Substituting $\alpha+2=n-\beta$ gives
$\beta(n-\beta)=n(n-1)/4$, a quadratic with roots $\bigl(n\pm\sqrt n\bigr)/2$.

Suppose $N$ is odd, so $f(u)\not\equiv e(u)$, $a$ is white and $b$ is black.
The pair $\{a,b\}$ has gap $0$ and is bichromatic.  The pair $\{a,u\}$ has gap
$e(u)$ and is bichromatic exactly when $u$ is black, i.e.\ when $e(u)$ is even;
the pair $\{b,u\}$ has gap $f(u)$ and is bichromatic exactly when $u$ is white,
i.e.\ when $e(u)$ is odd and hence $f(u)$ even.  Finally
$\operatorname{gap}(u,v)\equiv\varepsilon(u)+\varepsilon(v)+1$, which is even
exactly when $u$ and $v$ have different colours.  Hence the even gaps are
precisely the bichromatic pairs, of which there are $(1+\alpha)(1+\beta)$, and
$[0,N-1]$ contains $(N+1)/2$ even values.  With $\alpha+\beta=n-2$ this makes
$\alpha$ and $\beta$ the two roots of $t^2-(n-2)t+\tfrac{(n-2)(n-3)}4$, whose
discriminant is $n-2$.
\end{proof}

\begin{corollary}[Taylor's condition, recovered]\label{cor:detaylor}
If a Leech tree of order $n\ge3$ exists then $n$ is a perfect square (when
$n\equiv0,1\bmod4$) or $n-2$ is a perfect square (when $n\equiv2,3\bmod4$).
\end{corollary}

\begin{proof}
$N=\bin n2$ is even exactly for $n\equiv0,1\pmod4$.  In that case
Proposition~\ref{prop:desplit}(a) forces $\sqrt n\in\Z$, and in the other case
Proposition~\ref{prop:desplit}(b) forces $\sqrt{n-2}\in\Z$.
\end{proof}

Corollary~\ref{cor:detaylor} is Taylor's parity obstruction
(Lemma~\ref{lem:taylor}, \cite{Taylor1977}) in the two-anchor language.  The
point of Proposition~\ref{prop:desplit} is not the corollary, which is known,
but that it pins the split \emph{exactly}: at $n=25$ we have $N=300$, $V=23$,
$\beta(25-\beta)=150$, and therefore
\begin{equation}\label{eq:desplit25}
  (\alpha,\beta)\in\bigl\{(13,10),\,(8,15)\bigr\},
\end{equation}
so that a search may branch on two cases and enforce the counts as it goes.  At
$n=4$ the larger root $\beta=3$ exceeds $V=2$ and only $(\alpha,\beta)=(1,1)$
survives; for $n\ge9$ both roots are admissible.

\subsubsection*{Attachment, spine and balls}

\begin{lemma}[attachment]\label{lem:deattach}
Let $u\notin\{a,b\}$ with $h(u)>0$ and put $d=f(u)-e(u)$.  Then $T$ has a
vertex $w\notin\{a,b\}$ on $P$ with $p(w)=p(u)$, and its coordinates are
determined:
\[
  e(w)=\frac{N-d}2,\qquad f(w)=\frac{N+d}2,\qquad e(w)+f(w)=N .
\]
\end{lemma}

\begin{proof}
The vertex $q(u)$ lies on $P$ and satisfies $h(q(u))=0$ and $p(q(u))=p(u)$.  It
is not $a$: $a$ is a leaf, so its only neighbour lies on $P$, and $h(u)>0$
forces $q(u)$ to have a neighbour off $P$.  The same argument excludes $b$.  So
$w=q(u)$ is a non-anchor vertex with $h(w)=0$, and \eqref{eq:deph} with
$h(w)=0$ and $f(w)-e(w)=2p(w)-N=d$ gives the two values.
\end{proof}

Lemma~\ref{lem:deattach} turns every off-path vertex into a completely
determined pair of values that must be present, which is what makes it usable
as a prune.  It is a realisability condition rather than a search-order filter,
and it is the only condition of the search whose removal enlarges the solution
set in the ablation of Table~\ref{tab:deablate} --- see
Remark~\ref{rem:derelax} and the discussion in Section~\ref{sec:de:valid}.

\begin{lemma}[the spine is a Golomb ruler]\label{lem:despine}
The vertices of $T$ lying on $P$, together with $a$ and $b$, have pairwise
distinct distances, so their positions $p$ form a Golomb ruler of length $N$
with $k+2$ marks, where $k=\#\{u\notin\{a,b\}:h(u)=0\}$.  Hence
$G(k+2)\le N$, with $G$ as in Lemma~\ref{lem:golomb}.  At $n=25$, where
$N=300$, the published values $G(20)=283$ and $G(21)=333$ give $k\le18$, so at
least $5$ of the $23$ non-anchor vertices satisfy $h(u)>0$.
\end{lemma}

\begin{proof}
Distances between vertices of $P$ are differences of their positions, and all
distances in $T$ are distinct, so the positions form a Golomb ruler; its length
is $p(b)-p(a)=N$.  A Golomb ruler with $m$ marks has length at least $G(m)$, and
$G$ is strictly increasing, so $G(21)=333>300$ forces $k+2\le20$.
\end{proof}

We state plainly what Lemma~\ref{lem:despine} rests on.  The values $G(m)$ for
$m\le13$ were re-verified by exhaustive search for this project
(Section~\ref{sec:prelim}), but $G(20)=283$ and $G(21)=333$ are not; they are
taken from the published tables (OEIS A003022 \cite{OEIS}; see also
\cite{Shearer1990,Dollas1998}) and are the output of the distributed.net OGR
project~\cite{DistributedNetOGR}, a very large distributed computation which we
do not re-verify here and which is not part of the audit trail of
Section~\ref{sec:verification}.  Nothing else in this section depends on
Lemma~\ref{lem:despine}; it is used only as an orientation on how many vertices
must leave the spine, and it is not implemented as a prune.

\begin{lemma}[ball bound]\label{lem:deball}
Let $B$ be the set of vertices at distance at most $D$ from a fixed vertex of a
Leech tree.  Then all $\bin{|B|}2$ distances inside $B$ are distinct and at most
$2D$, so $\bin{|B|}2\le2D$.  Applied at $a$: for every $D\ge1$,
\[
  \#\{u\notin\{a,b\}: e(u)\ge N-D\}\ \le\ \frac{1+\sqrt{1+16D}}2-1 .
\]
\end{lemma}

\begin{proof}
For $x,y\in B$ the triangle inequality gives $d(x,y)\le2D$, and the distances
are distinct positive integers, so there are at most $2D$ of them.  Solving
$|B|(|B|-1)\le4D$ gives $|B|\le\bigl(1+\sqrt{1+16D}\bigr)/2$; at the centre $a$
the set $B$ consists of $a$ itself together with the vertices $u$ with
$d(a,u)\le D$, i.e.\ $e(u)\ge N-D$.
\end{proof}

\begin{remark}[cross-Sidon condition]\label{rem:decross}
The $\bin V2$ values \eqref{eq:degap} are pairwise distinct.  For untied pairs
this says that no four vertices satisfy $f(u)-f(u')=e(v')-e(v)\ne0$ with the
corresponding orderings, i.e.\ that the positive difference sets of $E$ and of
$F$ are almost disjoint.  This is a Sidon-type condition \emph{across} the two
sets rather than inside either of them, which is why the classical
Golomb-ruler bounds apply to the spine (Lemma~\ref{lem:despine}) but not to $E$
or $F$ separately.
\end{remark}

\subsection{A generating-function identity}\label{sec:de:lemmad}

Write $A(x)=\sum_i x^{e_i}$, $B(x)=\sum_i x^{f_i}$ and
$H(x)=\sum_i x^{e_i+f_i}$, the sums being over the $V$ non-anchor vertices.  For
an unordered pair $\{i,j\}$ let $m_{ij}\le M_{ij}$ be the two crossings
$f_i+e_j$ and $f_j+e_i$, so that
$m_{ij}+M_{ij}=(e_i+f_i)+(e_j+f_j)$ and, by Lemma~\ref{lem:degap},
$\operatorname{gap}(u_i,u_j)=m_{ij}+2h_c(i,j)$.

\begin{lemma}\label{lem:deD}
For a Leech tree of order $n$, as polynomials in $x$,
\begin{equation}\label{eq:deD}
  \bigl(A(x)+1\bigr)\bigl(B(x)+1\bigr)
  =\sum_{w=0}^{N-1}x^{w}\;+\;H(x)\;+\;\sum_{\{i,j\}}x^{M_{ij}}
   \;+\!\!\sum_{\substack{\{i,j\}\ \mathrm{tied}\\ h_c(i,j)>0}}\!\!
     \bigl(x^{m_{ij}}-x^{m_{ij}+2h_c(i,j)}\bigr).
\end{equation}
\end{lemma}

\begin{proof}
$AB=\sum_{i,j}x^{e_i+f_j}$; the diagonal terms give $H$ and each unordered pair
$\{i,j\}$ contributes $x^{m_{ij}}+x^{M_{ij}}$.  Writing
$x^{m_{ij}}=x^{m_{ij}+2h_c}+\bigl(x^{m_{ij}}-x^{m_{ij}+2h_c}\bigr)$ and using
$m_{ij}+2h_c=\operatorname{gap}(u_i,u_j)$ turns $AB$ into
$H+\sum_{\{i,j\}}\bigl(x^{\operatorname{gap}(u_i,u_j)}+x^{M_{ij}}\bigr)$ plus
the tie corrections, which vanish unless $h_c>0$.  Now
$(A+1)(B+1)=1+A+B+AB$, and by Theorem~\ref{thm:denormal}
$1+A+B+\sum_{\{i,j\}}x^{\operatorname{gap}(u_i,u_j)}=\sum_{w=0}^{N-1}x^w$.
\end{proof}

Counting terms on both sides of \eqref{eq:deD} gives
$(V+1)^2=N+V+\bin V2$, which is an identity; so the excess of the sumset
$\bigl(E\cup\{0\}\bigr)+\bigl(F\cup\{0\}\bigr)$ over the interval $[0,N-1]$ is
exactly $V+\bin V2$ terms, accounted for by $H$ and by the maximal crossings.

\begin{corollary}[covering]\label{cor:decover}
Put $E'=E\cup\{0\}$ and $F'=F\cup\{0\}$.  Every $w\in[0,N-1]$ is representable
as $e+f$ with $e\in E'$, $f\in F'$, except possibly those $w$ that are the gap
of a tied pair with $h_c>0$.
\end{corollary}

\begin{proof}
Every term of \eqref{eq:deD} other than the last sum is non-negative, and the
last sum subtracts $1$ from the coefficient of $x^{\operatorname{gap}(u_i,u_j)}$
for each tied pair with $h_c>0$; those exponents are pairwise distinct because
the gaps are, so at most $1$ is subtracted at any single $w$.
\end{proof}

\begin{remark}[the corollary is a tautology; the multiplicities are not]\label{rem:decaveat}
Corollary~\ref{cor:decover} is not a constraint.  If $w$ is the gap of an untied
pair then $w=f_u+e_v$ by definition, so it is representable by construction; if
$w\in E$ then $w=w+0$, and likewise for $F$ and for $w=0$.  The same remains
true in a partial state of the search, because both coordinates involved are
smaller than $w$ and hence already assigned.  We implemented the corollary as a
prune anyway, in order to measure it: at $n=9$ it changed the node count by
exactly zero ($24{,}335$ nodes with and without it; row \tok{--no-cover} of
Table~\ref{tab:deablate}).  That is the expected outcome and we record it so
that the corollary is not mistaken for a live constraint.

What Lemma~\ref{lem:deD} does say, and what is not a tautology, is a statement
about \emph{multiplicities}: for $w\in[0,N-1]$ the number of representations of
$w$ as $e+f$ with $e\in E'$, $f\in F'$, counted with multiplicity, equals
\[
  1+\#\{i:e_i+f_i=w\}+\#\{\{i,j\}:M_{ij}=w\}
\]
up to the tie correction.  Since $e_i+f_i=N-2h_i$ and
$M_{ij}=N+|p_i-p_j|-h_i-h_j$, the sumset multiplicities read off the heights and
the maximal crossings.  Counting the representations that reach $N$ gives, with
$k$ the number of on-path vertices as in Lemma~\ref{lem:despine},
\[
  \#\{(e,f)\in E\times F: e+f\ge N\}
  \;=\;k+\#\bigl\{\{i,j\}:|p_i-p_j|\ge h_i+h_j\bigr\}.
\]
\end{remark}

\subsection{The gap-order search}\label{sec:de:search}

The normal form suggests a search order that the tree structure does not: since
every element of $E$ and of $F$ is itself a gap, one may process the target
values $g=1,2,\dots,N-1$ in increasing order, maintaining the invariant that
every gap smaller than $g$ has been realised and every coordinate not yet
assigned is at least $g$.  Vertices are created with only one of their two
coordinates known; the partner is supplied later, at its own value of $g$.

\begin{lemma}[determination]\label{lem:dedet}
Suppose that at step $g\ge1$ every value in $[0,g-1]$ is realised by a pair
whose gap is already determined, and that every coordinate not yet assigned is
$\ge g$.  Let $u,v$ be non-anchor vertices, let $X=f(u)+e(v)$ and $Y=f(v)+e(u)$
be their crossings, and let $X^-,Y^-$ be the values obtained from $X,Y$ by
replacing every unassigned coordinate by $g$.  Then
\begin{enumerate}
\item[(a)] if both crossings are known, $\operatorname{gap}(u,v)=\min(X,Y)$ when
$X\ne Y$, and $\operatorname{gap}(u,v)=X+2h_c$ with
$0\le h_c\le\min(h_u,h_v)$ when $X=Y$;
\item[(b)] if exactly one crossing, say $X$, is known and $X<Y^-$, then
$\operatorname{gap}(u,v)=X$;
\item[(c)] in every other case $\operatorname{gap}(u,v)\ge g+1>g$.
\end{enumerate}
\end{lemma}

\begin{proof}
Part (a) is Lemma~\ref{lem:degap}.  A crossing with an unassigned coordinate is
at least $1+g$, the other summand being at least $1$ by
Lemma~\ref{lem:decoord}(i), so $X^-,Y^-\ge g+1$ whenever the corresponding
crossing is unknown.  In case (b), $X<Y^-\le Y$, so the crossings are unequal,
hence $h_c=0$ and the gap is $\min(X,Y)=X$.  In case (c) either both crossings
are unknown, and then $\operatorname{gap}(u,v)\ge\min(X^-,Y^-)\ge g+1$, or
exactly one is known with $X\ge Y^-\ge g+1$, and then
$\operatorname{gap}(u,v)\ge\min(X,Y)\ge\min(X,Y^-)\ge g+1$.
\end{proof}

\begin{corollary}[two options at each step]\label{cor:debranch}
Assume the hypotheses of Lemma~\ref{lem:dedet}, and suppose that cases (a) and
(b) have been applied to every pair at the current value of $g$.  If $g$ is
still unrealised then either $g\in E\cup F$, i.e.\ $g$ is assigned as a
coordinate of an existing half-known vertex or of a new vertex, or $g$ resolves
a previously deferred tied pair, i.e.\ $g=X+2h_c$ for a pair with $X=Y$ known
and $h_c>0$.
\end{corollary}

\begin{proof}
The value $g$ is realised by some pair.  That pair is not $\{a,b\}$, whose gap
is $0$.  If it is $\{a,u\}$ or $\{b,u\}$ then $g\in E\cup F$.  Otherwise it is a
pair $\{u,v\}$ of non-anchor vertices.  Case (c) of Lemma~\ref{lem:dedet} would
give $\operatorname{gap}(u,v)>g$, and cases (b) and the untied part of (a) have
by assumption already been applied, so the only remaining possibility is a tied
pair, whose gap is $X+2h_c$ with $h_c>0$ still to be chosen.
\end{proof}

Corollary~\ref{cor:debranch} makes the search tree shallow in a precise sense:
every other gap is forced, so the only decisions are the $2V$ coordinate
assignments together with the resolutions of deferred ties.

\paragraph{The rescan, and why it is not optional.}
The threshold in case (b) of Lemma~\ref{lem:dedet} is $X<Y^-$, and $Y^-$ grows
with $g$.  A pair whose known crossing was too large to be conclusive at one
step therefore becomes determined at a later step \emph{without any new
assignment}, as soon as $g$ has risen far enough.  The engine must consequently
rescan all still-undetermined pairs at every value of $g$, mark those that have
become determined, and only then decide whether $g$ itself is still unrealised;
this is also what makes the hypothesis of Corollary~\ref{cor:debranch} true when
the branching is reached.  The comparison must be against the computed bound
$Y^-$ and never against a constant.  The invariant is that an unassigned
coordinate is $\ge g$ while the step for $g$ is being decided, and tightens to
$>g$ only once the value $g$ has been assigned; a fixed threshold derived from
the tighter form --- testing $X\le g+1$ rather than $X<Y^-$ --- would be unsound
at the moment of decision, since the pair could then be tied with $X=Y=g+1$ and
a gap larger than $X$.  Omitting the rescan does not merely slow the search
down: the gap $X$ is then never registered, its value can be handed out a second
time, and the search reports \emph{no solution} where solutions exist.  This was
one of the four bugs discussed in Section~\ref{sec:de:valid}.

\paragraph{The prunes.}
The engine applies the following conditions, each of which can be switched off
individually; the flag names are those of the ablation in
Table~\ref{tab:deablate}.
\begin{itemize}\itemsep2pt
\item \emph{partner feasibility} (always on): a half-known vertex must still
admit an untaken partner value $\ge g$ of the parity required by
Lemma~\ref{lem:decoord}(iii) and with $e+f\le N$.
\item \emph{deferred feasibility} (\tok{--no-\bk deferred-\bk check}): a deferred
tied pair must still admit a resolution, i.e.\ an untaken value of the right
parity in its window $[X,\,X+2\min(h_u,h_v)]$ that is $\ge g$.
\item \emph{partner matching} (\tok{--no-hall}): distinct half-known vertices
need distinct partner values, so a bipartite matching between half-known
vertices and admissible values must exist (Hall's condition, tested by
augmenting paths).
\item \emph{lower-bound Hall} (\tok{--no-lb}): every untaken value in $[g,N)$
must be supplied either by one of the remaining coordinate assignments or by a
pair whose gap lower bound is small enough; a counting test over sorted lower
bounds.
\item \emph{upper-bound Hall} (\tok{--no-ub}): the dual test, in which every
untaken value must be supplied by something whose upper bound is large enough.
\item \emph{attachment} (\tok{--no-attach}): Lemma~\ref{lem:deattach}, i.e.\
every vertex with $h>0$ forces a determined on-path vertex at the same $p$,
whose two values must be present or still free.
\item \emph{meeting vertex} (\tok{--no-hc}): when a tie is resolved with
$h_c>0$, the meeting point is itself a vertex of $T$ at height $h_c$ in the same
$p$-class, so its two forced coordinates must be present or still free.
\item \emph{ultrametric class test} (\tok{--no-ultra}): within one $p$-class the
meeting heights form a max-ultrametric --- among any three, the smallest occurs
at least twice --- and, since two vertices of a class cannot share a height, a
meeting at height $h_c>0$ identifies a \emph{unique} class vertex, which must
then be an ancestor of both endpoints.
\item \emph{covering} (\tok{--no-cover}): Corollary~\ref{cor:decover}; see
Remark~\ref{rem:decaveat}, where it is shown to be worth nothing.
\item \emph{parity} (\tok{--no-parity}): the running counts
$(\alpha,\beta)$ must stay inside an admissible split of
Proposition~\ref{prop:desplit}.
\end{itemize}

\paragraph{Symmetry breaking.}
By Theorem~\ref{thm:endpoint} the value $N-1$ is attained by a pair containing
an anchor, i.e.\ $\operatorname{gap}=1$ is a coordinate, and exchanging $a$ with
$b$ exchanges $E$ with $F$.  One may therefore assume $1\in E$.  This is the
only symmetry of the normal form, and it halves the search exactly: at $n=9$,
$24{,}335$ nodes with it against $48{,}669=2\cdot24{,}335-1$ without
(Table~\ref{tab:desmall}).

\subsection{Validation}\label{sec:de:valid}

The engine exists in two implementations.  The Python reference is
\file{double\_\bk end\_\bk search.py}, and
\file{double\_\bk end\_\bk search.cpp} is a single-file C++17 port of it; both
are in the archive directory \file{theory-lab/\bk double\_\bk end/}
(Section~\ref{sec:data}).  The code was written with AI assistance
(Section~\ref{sec:ai}) and was treated as untrusted; what follows is what was
actually done to it.

\paragraph{Differential test against an independent enumerator.}
\file{brute\_\bk double\_\bk end.py} enumerates the pure relaxation of
Remark~\ref{rem:derelax} in a completely different order --- over vertices, in
increasing order of $e$ --- and checks the defining conditions directly, with no
attachment, meeting-vertex or ultrametric conditions and no gap-order
reasoning.  With every optional prune of the gap-order engine switched off, the
two agree exactly, as sets of coordinate systems $\{(e_i,f_i)\}$, for
$n=3,4,6$: $2$, $4$ and $12$ of them respectively.  (The comparison is made on
coordinate systems because the engine reports each admissible choice of tie
heights separately, so its raw solution count can be larger: at $n=6$ the
twelve coordinate systems appear as sixteen reports.)  The twelve at $n=6$ are
solutions of the relaxation only; with the prunes on the engine returns $2$
(the single Leech tree of order $6$ in its two anchor orientations, or $1$ after
symmetry breaking), and every one of the twelve that rebuilds as a genuine tree
is among them.

\paragraph{Genuine trees, rebuilt and checked.}
Every reported solution is converted back into an explicit weighted tree --- the
on-path chain ordered by $p$, the off-path branches grouped by $p$ and attached
at the recorded meeting heights --- and all its pairwise distances are then
recomputed by breadth-first search and compared with $[1,N]$.  With all prunes
on and symmetry breaking off, the engine returns exactly the known Leech trees
of orders $3$, $4$ and $6$, each once per anchor orientation, and all of them
pass the rebuild check.

\paragraph{Independent brute force.}
\file{brute\_\bk leech\_\bk small.py} is a from-scratch enumerator of Leech trees
themselves, sharing no code with any engine in this project: it generates all
tree shapes, searches weights, and canonicalises the weighted trees.  It gives
the complete counts
\[
  1,\;1,\;2,\;0,\;1,\;0,\;0,\;0\qquad(n=2,\dots,9)
\]
up to isomorphism, in agreement with the literature: the five known Leech trees
are those of Table~\ref{tab:known}~\cite{Leech1975}, the orders $n=5,7,8$ are
excluded by Taylor's condition~\cite{Taylor1977} and $n=9$ by the searches
of~\cite{SzekelyWangZhang2005,Calhoun2007}.  The
gap-order engine reproduces exactly these counts
(Table~\ref{tab:desmall}); the three orders with count $0$ and no search tree at
all, $n=5,7,8$, are the orders excluded outright by
Proposition~\ref{prop:desplit}.  The oracle is archived under the same
discipline as the engines: both the enumerator and the counts it produced are
frozen and hash-tagged (rows~6 and~7 of Table~\ref{tab:dehash}), so the
regression fixture the engine is tested against is itself fixed.

\paragraph{Trace-identical port.}
The C++ port was accepted on a stronger criterion than agreeing node counts.  It
emits, under \tok{--trace}, one line per recursive call carrying the full state
$(g,\ \text{number of vertices},\ e[\,],\ f[\,],\ |\text{deferred}|)$, and the
same trace is produced from the Python by instrumenting its recursion.  For all
eight configurations $n\in\{3,4,6,9\}$ with symmetry breaking on and off, the
two traces are byte-identical: the C++ visits the same nodes, in the same order,
in the same state.  The solution sets, compared as sorted sets of $(e_i,f_i)$
pairs, are identical as well.  That comparison was made against the first frozen
Python reference (row~1 of Table~\ref{tab:dehash}); the port was then re-checked
against the second freeze (row~2), which is the archived reference and differs
from the first only by the stable-deferred restore described below, and
reproduces its numbers.

\paragraph{Sharding.}
For cluster runs the search is split by \tok{--split LEVEL K M}, which numbers,
in depth-first order, the branches taken from nodes at recursion depth
\tok{LEVEL} and explores those congruent to $K$ modulo $M$.  That the shards
partition the run was verified by exhaustion rather than asserted: with $A$ the
number of nodes at depth $\le\tok{LEVEL}$ and $B$ the number of deeper nodes,
the shards must together visit $M\!\cdot\!A+B$ nodes, and the multiset of visited
states must be $M$ copies of the shallow part plus one copy of each deep node.
Both were checked by dumping the complete trace of every shard and of the
unsharded run and comparing multisets, for six configurations at $n=6$ and
$n=9$; the solutions of the shards likewise union to the unsharded solutions
with no solution appearing twice.

\paragraph{Prune ablation.}
Table~\ref{tab:deablate} gives the cost of each condition at $n=9$.  Eight of
the nine leave the solution set literally unchanged.  The attachment condition
does not, and this is the expected direction: it is a realisability condition,
not a search-order filter, so switching it off enlarges the relaxation.  The
soundness test that matters holds in both directions --- no baseline solution
is ever lost, and each of the $18$ extra solutions at $n=9$ (and $5$ at $n=6$)
is rejected by the rebuild check, every one of them with the message
\tok{no on-path attachment vertex at p=\dots}, which is exactly the condition of
Lemma~\ref{lem:deattach}.  So no prune drops a genuine tree, and the attachment
condition removes only non-trees.

\begin{table}[ht]
\centering
\caption{Prune ablation at $n=9$ (symmetry breaking on; every run exhausted).
``Solutions'' counts solutions of the relaxation, not trees: the baseline is
empty, and the $18$ solutions of the \tok{--no-attach} row are all rejected by
the rebuild-and-BFS check.  Every row except \tok{--no-hc} was re-run for this
paper; that row is quoted from the port record.}
\label{tab:deablate}
\begin{tabular}{@{}llrr@{}}
\toprule
condition disabled & flag & nodes & solutions\\
\midrule
(none)                 & ---                    & $24{,}335$    & $0$\\
partner matching       & \tok{--no-hall}        & $24{,}335$    & $0$\\
lower-bound Hall       & \tok{--no-lb}          & $24{,}335$    & $0$\\
covering               & \tok{--no-cover}       & $24{,}335$    & $0$\\
parity split           & \tok{--no-parity}      & $24{,}450$    & $0$\\
deferred feasibility   & \tok{--no-\bk deferred-\bk check} & $25{,}997$ & $0$\\
upper-bound Hall       & \tok{--no-ub}          & $28{,}338$    & $0$\\
attachment             & \tok{--no-attach}      & $105{,}230$   & $18$\\
ultrametric class test & \tok{--no-ultra}       & $1{,}462{,}492$ & $0$\\
meeting vertex         & \tok{--no-hc}          & $11{,}942{,}021$ & $0$\\
\bottomrule
\end{tabular}
\end{table}

The single largest structural gain came from the last condition added to the
ultrametric test: within one $p$-class two vertices cannot share a height, since
equal $p$ and equal $h$ force equal $(e,f)$, so a meeting at height $h_c>0$
identifies a unique class vertex, which must be an ancestor of both endpoints.
In the development record that condition alone took $n=9$ from $62{,}612$ to
$24{,}335$ nodes; the intermediate build is not part of the archive, so unlike
every row of Table~\ref{tab:deablate} this figure was not re-measured.

\paragraph{Four bugs, and what was discarded.}
Four genuine errors were found in the course of this validation, and we record
them because each of them produced a wrong answer of the dangerous kind --- a
spurious ``no solution'' --- and because every number measured before the
corresponding fix was discarded rather than reconciled.
\begin{enumerate}\itemsep2pt
\item An unsound meeting-vertex test.  The first version of the meeting-vertex
condition (\tok{--no-hc}) required the meeting vertex of a tied pair to \emph{already exist} in
the partial state, rather than allowing it to be created later, and so rejected
valid states.
\item An unstable index in the deferred-pair loop.  The loop over deferred tied
pairs deleted entries from the list it was iterating over, so live indices
shifted and some deferred pairs were skipped.
\item The missing rescan.  As explained in Section~\ref{sec:de:search}, an
undetermined pair becomes determined as $g$ grows even when nothing is assigned.
The original engine tested for determination only at the moment of assignment,
and therefore lost gaps.
\item A non-state-neutral undo.  The undo of a resolved deferred pair restored
the entry at the \emph{end} of the deferred list instead of at the index it was
removed from.  Node counts and solution sets are invariant under that
permutation, so unsharded runs were unaffected --- but depth-first order is not,
and \tok{--split} numbers branches in depth-first order, so shards silently
drifted.  At $n=9$ with $\tok{LEVEL}=10$ and $M=4$ the shards visited $69{,}733$
nodes where the partition identity requires $70{,}201$: $468$ nodes were lost.
The restore now happens at the recorded index, which makes the recursion
state-neutral; this changes neither node counts nor solution sets at
$n=3,4,6,9$.
\end{enumerate}
A fifth error, in the C++ port only, was a scratch candidate list held as a
member and clobbered by recursion; it was located by diffing the C++ trace
against the Python trace and identifying the first divergent state.  It is
mentioned here because it is the concrete reason the trace-level comparison was
adopted in place of comparing node counts.

\paragraph{What was re-run for this paper.}
On the author's machine, both implementations were re-run and agree:
$3/5$, $8/15$, $44/87$ and $24{,}335/48{,}669$ nodes at $n=3,4,6,9$ with and
without symmetry breaking, with $1/2$, $2/4$, $1/2$ and $0/0$ solutions; the
differential test against the independent enumerator was re-run at $n=3,4,6$;
the ablation of Table~\ref{tab:deablate} was re-run except for the
\tok{--no-hc} row; the rebuild check was re-run on the $18$ extra solutions
of the \tok{--no-attach} row, all $18$ of which are rejected; the digests of
Table~\ref{tab:dehash} were recomputed from the archived files; and the
closed forms of Proposition~\ref{prop:desplit} were checked against the
admissible splits computed by the engine for every $n\le38$ and against the
coordinates of all five known Leech trees.
Table~\ref{tab:dehash} gives the SHA-256 digests of the archived sources.  The
Python reference was frozen twice on the same day; the two freezes differ only
in the stable-deferred restore of item~4 below, and produce the same eight node
counts and the same solution sets at $n=3,4,6,9$.  Both digests are recorded
because the trace comparison of the port was made against the first and
re-checked against the second.  Each archived file is named by the first
sixteen hexadecimal digits of its own digest, so that no name can come to refer
to two different files.

\begin{table}[ht]
\centering
\caption{SHA-256 digests of the archived sources of this section.  Row~1 is the
first frozen Python reference, against which the trace-identical comparison of
the C++ port was made; row~2 is the current reference, which adds the
stable-deferred restore and against which the port was re-checked.  Row~3 is
the C++ port, rows~4 and~5 the independent enumerator and the differential
harness, rows~6 and~7 the independent brute-force oracle and the counts it
produced, and row~8 the derivation these results follow.}
\label{tab:dehash}
\footnotesize
\setlength{\tabcolsep}{4pt}
\begin{tabular}{@{}rl@{}}
\toprule
\# & object and digest\\
\midrule
1 & first frozen Python reference (digest only; superseded)\\
  & \texttt{\scriptsize 934b24162db8c0c2b2d88e10d4b0b7d0adc78eb70da42556bb563ceea1884eb3}\\
2 & \file{frozen/double\_\bk end\_\bk search\_\bk 4916cae4bfc30da7.py}\\
  & \texttt{\scriptsize 4916cae4bfc30da753ce74532a64c74f0b24366f8ad4222122488890db2b5b1d}\\
3 & \file{frozen/double\_\bk end\_\bk search\_\bk 4eaddc2f3c1e28ce.cpp}\\
  & \texttt{\scriptsize 4eaddc2f3c1e28ce8b9f07a1879109e72b4a4b2c1f3d7c792ef848a73f97cf2b}\\
4 & \file{frozen/brute\_\bk double\_\bk end\_\bk 8be4ae51d96b0469.py}\\
  & \texttt{\scriptsize 8be4ae51d96b046936a0820d70103d7cbf5f5642cc427a85bad4e0616c319a2c}\\
5 & \file{frozen/differential\_\bk test\_\bk d82097986ad0afc7.py}\\
  & \texttt{\scriptsize d82097986ad0afc75f290d774adfca6991ba22602b66a5948188ef2e968d294b}\\
6 & \file{frozen/brute\_\bk leech\_\bk small\_\bk 93cafc42f9c27c0a.py}\\
  & \texttt{\scriptsize 93cafc42f9c27c0a6e3328f41e4aa09a3ed1f47f4738eeff18c6061c268028c4}\\
7 & \file{frozen/ground\_\bk truth\_\bk small\_\bk 34ae25138560c313.json}\\
  & \texttt{\scriptsize 34ae25138560c313d0e6f66214a05a462dc2a424abc46e55ffd26b21c1911ea0}\\
8 & \file{THEORY.md}, the derivation these results follow\\
  & \texttt{\scriptsize 58ffe0ea6c8b4aec57fd10dda02fa297942a7e1c65fb306418871971d1e43b7b}\\
\bottomrule
\end{tabular}
\end{table}

\begin{table}[ht]
\centering
\caption{The gap-order engine on the orders it has exhausted.  $(\alpha,\beta)$
are the admissible parity splits of Proposition~\ref{prop:desplit}; when there
are none the search tree is empty and no node is expanded.  ``Solutions'' are
solutions of the normal form, each of which was rebuilt as an explicit tree and
verified by breadth-first search; without symmetry breaking each tree appears
once per anchor orientation.  The last column is the independent brute-force
count of Leech trees up to isomorphism; that oracle covers $n\le9$, and the
non-existence at $n=11$ is known from the literature instead.  The $n=11$ run
was made with symmetry breaking only, so the two cells for the free search are
not measured.}
\label{tab:desmall}
\begin{tabular}{@{}rrrlrrrrr@{}}
\toprule
$n$ & $N$ & $V$ & $(\alpha,\beta)$ & \multicolumn{2}{c}{nodes} &
\multicolumn{2}{c}{solutions} & brute force\\
 & & & & sym.\ & free & sym.\ & free & \\
\midrule
$3$ & $3$  & $1$ & $(0,1),(1,0)$ & $3$        & $5$        & $1$ & $2$ & $1$\\
$4$ & $6$  & $2$ & $(1,1)$       & $8$        & $15$       & $2$ & $4$ & $2$\\
$5$ & $10$ & $3$ & none          & $0$        & $0$        & $0$ & $0$ & $0$\\
$6$ & $15$ & $4$ & $(1,3),(3,1)$ & $44$       & $87$       & $1$ & $2$ & $1$\\
$7$ & $21$ & $5$ & none          & $0$        & $0$        & $0$ & $0$ & $0$\\
$8$ & $28$ & $6$ & none          & $0$        & $0$        & $0$ & $0$ & $0$\\
$9$ & $36$ & $7$ & $(1,6),(4,3)$ & $24{,}335$ & $48{,}669$ & $0$ & $0$ & $0$\\
$10$ & $45$ & $8$ & none          & $0$        & $0$        & $0$ & $0$ & ---\\
$11$ & $55$ & $9$ & $(3,6),(6,3)$ & $3{,}338{,}887{,}905$ & --- & $0$ & --- & ---\\
\bottomrule
\end{tabular}
\end{table}

\subsection{Status, and what is not claimed}\label{sec:de:status}

\paragraph{Exhausted orders.}
The gap-order engine has been run to exhaustion for $n\le11$, with the results
of Table~\ref{tab:desmall}.  At $n=9$ it expands $24{,}335$ nodes with symmetry
breaking and $48{,}669$ without; at $n=11$, $3{,}338{,}887{,}905$ nodes with
symmetry breaking.  Neither returns a solution, and the orders $n=5,7,8,10$ are
disposed of by Proposition~\ref{prop:desplit} without expanding a node.  In the
terminology we use for finite computations, this is a
\tok{CERTIFIED FINITE IMPOSSIBLE} for the range $3\le n\le11$: within that
range, and only within it, the computation establishes that there is no Leech
tree beyond the four the search itself returns (one of order $3$, two of order
$4$, one of order $6$).  It agrees with the independent brute force for
$n\le9$, which is as far as that oracle goes, and with the literature at
$n=11$.  It is not, and must not be read as, evidence about any larger order.

\paragraph{Order 11.}
The $n=11$ search was run as a single process on a $128$-core workstation and is
\tok{EXHAUSTED} with no solution after $3{,}338{,}887{,}905$ nodes, with
symmetry breaking on.  It took about nine hours of wall clock on a machine whose
load average stayed between $200$ and $400$ throughout, so the CPU time was far
below the elapsed time and the wall figure measures the machine, not the
program.  This is an independent re-derivation of the non-existence of a Leech
tree of order $11$, by an algorithm and a normal form entirely different from
those used elsewhere in this paper.

\paragraph{Orders $n\ge16$ are out of reach for this engine.}
With $n=11$ measured, the node counts at the admissible orders are
\[
  3,\quad 8,\quad 44,\quad 24{,}335,\quad 3{,}338{,}887{,}905
  \qquad\text{at } n=3,4,6,9,11,
\]
that is at $V=1,2,4,7,9$.  The growth per unit of $V$ between consecutive
measured orders is about $2.7$, $2.3$, $8.2$, and then $370$.  The rate is not
constant; it is accelerating.  Continuing at the last measured rate would put
$n=16$ near $10^{22}$ nodes, of order $10^{12}$ CPU-hours, and $n=18$ far beyond
that.  We state plainly that $n=16$ and $n=18$ are out of reach for this engine,
and equally plainly that the extrapolation which says so is unreliable, because
a rate that has moved from $2.3$ to $370$ over the measured range carries no
guarantee of staying at $370$.  Both readings point the same way, so we make no
reachability claim beyond what has been measured.  This corrects an earlier
estimate twice over: a fit to $n\le9$ predicted some $2\times10^6$ nodes at
$n=11$ and was first refuted by the C++ measurements, which had not terminated
after $2.5\times10^7$ nodes; the measured sequence above now settles it.  A
sampled measurement at $n=16$ was planned and has been cancelled as pointless.

What the engine is for, then, is the independent re-derivation of the known
non-existence results up to $n=11$ and the restricted order-$25$ search below,
not an attack on $n=16$ or $n=18$.

\paragraph{An independent route to the order-25 certificates.}
The two-anchor language connects to Section~\ref{sec:crowding} in one clean way.

\begin{proposition}\label{prop:des}
Let $T$ be a Leech tree of order $n$ with anchors named as in
Theorem~\ref{thm:endpoint}, and put
$s=\min\bigl(F\cup\{\operatorname{gap}(u,v):u,v\notin\{a,b\}\}\bigr)$.  Then
$\diam(T-a)=N-s$, and $1,2,\dots,s-1$ all lie in $E$.  Consequently, at $n=25$,
\[
  \diam(T-a)\le284
  \iff s\ge16
  \iff \{1,2,\dots,15\}\subseteq E .
\]
\end{proposition}

\begin{proof}
The pairs of $T-a$ are the pairs $\{b,u\}$ and the pairs $\{u,v\}$ with
$u,v\notin\{a,b\}$, whose gaps are the elements of $F$ and the values
\eqref{eq:degap}; the largest distance is $N$ minus the smallest gap, so
$\diam(T-a)=N-s$.  Every value in $[1,s-1]$ is realised, by
Theorem~\ref{thm:denormal}, and not by a pair avoiding $a$, and not by
$\{a,b\}$, so it lies in $E$.  Conversely if $\{1,\dots,15\}\subseteq E$ then,
the $N$ gaps being distinct, no element of $F$ and no pair gap is $\le15$, so
$s\ge16$.  With $N=300$, $\diam(T-a)=300-s\le284$ is $s\ge16$.
\end{proof}

So the hypothesis $\diam(T-a)\le284$, which is what
Theorem~\ref{thm:crowd280} and the four finite certificates of
Theorem~\ref{thm:crowd283} together exclude, is in the two-anchor language the
completely explicit restriction ``the first fifteen gaps are all $e$-values''.
The engine implements it directly (\tok{--force-e-upto 15}).  That search is
running, as $128$ shards split at recursion level $44$ --- above the
forced-chain level $15$, for the reason given below --- of which four are on a
workstation and the rest are queued on the cluster; nothing has come back from
it, and we claim nothing from it here.  If it is exhausted and returns nothing,
it
re-derives the conclusion $\diam(T-a)\ge285$ --- that is, the content of
Theorems~\ref{thm:crowd280} and~\ref{thm:crowd283} for the endpoint $a$, which
is the only leaf Corollary~\ref{cor:reduction} uses, covering the four finite
certificates $\delta=281,\dots,284$ and every smaller $\delta$ at once --- by a
different algorithm, in a different normal form, with no shared code.  That is
the two-implementation standard this project applies to every load-bearing
computation, and it is the most valuable run the engine could make.  We stress
two things.  First,
\tok{--force-e-upto} is a \emph{case restriction}, not a prune: it deliberately
loses solutions (at $n=4$ the solution count drops from $2$ to $0$ at $K=3$, and
at $n=6$ from $1$ to $0$ at $K=5$), so a zero result under it certifies only the
restricted statement and says nothing about $\delta\ge285$.  Second, the
restricted search is locked by its first fifteen steps but is otherwise a full
order-$25$ search, and its size has not been measured.
\paragraph{How far the forcing goes, and how we know.}
Under \tok{--force-e-upto S} each of the first $S$ steps has exactly one legal
option: every vertex created so far already carries its $e$-coordinate, and
$f$-options are forbidden while $g\le S$, so the only candidate is a new vertex
with $e=g$.  The search is therefore a forced chain down to level $S$, and a
\tok{--split} taken at any level below $S$ numbers a single branch, putting the
whole search into shard $0$ and leaving every other shard with the prefix alone.
Anyone sharding a restricted run should split deeper than $S$, or not at all.

The forcing stops there.  For $n=25$ with $S=15$, shard $1$ of $1000$ is
\tok{EXHAUSTED} after $14$ nodes when split at level $13$ and after $15$ nodes
when split at level $14$ --- the prefix only, so those levels hold a single node
--- while at level $15$ and beyond the same shard is still running after
millions of nodes.  The first free decision is at gap $S+1$, and there is no
forcing past it.

We record how this was established, because our first attempt to establish it
was wrong.  The per-depth histogram printed by \tok{--debug} after a run that
stopped at a node limit does \emph{not} show that a level is forced: depth-first
search descends before it backtracks, so a shallow level shows a count of $1$
merely because its remaining siblings have not been reached yet.  An early
reading of such a histogram suggested a $24$-step forced chain, with $E$
continuing $18,36,54,\dots$ and $F$ continuing $16,32,48,\dots$; that is the
first depth-first path and nothing more, and both the shard test above and an
exhaustive small-order run refute it.  The exhaustive check is $n=9$ with
\tok{--force-e-upto 5}, where depths $0$ to $5$ hold one node each and the first
branching is at depth $6$.  Only an exhaustive run, or a shard test of the kind
above, decides forcedness.

That re-derivation was attempted and then abandoned, and we report the
measurement that ended it, because it is the honest cost of the method.

The engine keeps no checkpoint, so a shard stopped by the wall clock certifies
nothing; a first attempt with a four-hour limit therefore returned nothing at
all from $128$ shards.  A second attempt with $480$ shards and a $23$-hour
limit gave the rate: $100$ shards explored $1.96\times10^{11}$ nodes in four
hours, that is about $1.4\times10^{5}$ nodes per second per core, roughly an
order of magnitude below the rate at $n\le11$ because the prunes of
Section~\ref{sec:de:search} cost $O(N+V^2)$ per node and $N=300$ here.  We then
sized the search directly, by drawing $200$ shards from a $20{,}000$-way split
at recursion level $47$ and giving each a two-hour limit.  Every one of the
$100$ that ran reached the wall without exhausting its shard.  At the measured
rate each of those shards had passed $9.8\times10^{8}$ nodes, so

\[
  \text{total} \;>\; 20{,}000 \times 9.8\times10^{8} \;=\; 1.96\times10^{13}
  \text{ nodes} \;\approx\; 4\times10^{4}\text{ CPU-hours},
\]

and this is a floor, not an estimate: the true figure is whatever it takes to
finish shards none of which finished.  For comparison, the $\delta=284$
certificate of Theorem~\ref{thm:crowd283} is $1.64\times10^{9}$ structures, and
$\delta=285$ was costed at about $10^4$ CPU-hours.

The conclusion is negative and worth stating plainly: at order $25$ the
two-anchor search is at least an order of magnitude more expensive than the
$\delta$-indexed programme it would have checked, so the independent
re-derivation of Theorem~\ref{thm:crowd283} was not carried out and nothing in
this section rests on it.  The parameter-free normal form buys generality at
the cost of the structure that makes the $\delta$-indexed search cheap: fixing
$\delta$ fixes the top block, and the crowding of Section~\ref{sec:crowdineq}
then prunes it, whereas here every admissible $\delta$ is explored at once.
What the two-anchor engine does deliver is Theorem~\ref{thm:denormal} and
its consequences, and the independent re-derivation of non-existence for
$n\le11$.

Conversely, the still-open range $285\le\delta\le298$ of
Corollary~\ref{cor:reduction} corresponds to $s\le15$, which is a \emph{weak}
restriction on the two-anchor search: the existing order-$25$ certificates
therefore do very little to help a full order-$25$ run in this normal form.  The
two approaches are complementary only in the direction described above.
